\section{Proposed Solution}
\label{sec:proposed}

GRASP provides a connected workflow for reviewing company BRSR reports. The proposed solution brings text segments, reported figures, principle assignments, and company scores into one system. An analyst can examine the evidence behind a review flag and the calculation behind a score. The overall architecture is shown in Figure~\ref{fig:architecture}; the purpose of the design is to make disclosure review organised, traceable, and easier to interpret.

The first design choice addresses the difficulty of connecting claims with evidence, described in Section~\ref{sec:introduction}. Manual review requires an analyst to search for a statement, locate the relevant figures, and check the reporting context. GRASP stores this information by company, report, year, and reporting item. Language checks retain text segments marked for review. Numerical checks create records containing a claim and the relevant reported figures, together with available assurance information about external verification. These records give an analyst the information needed to examine a possible mismatch.

The second design choice uses the regulatory structure to organise training data. Principle tags, reporting questions, and numbered principle notes provide rules for assigning labels to text segments. These automatically assigned labels are called \emph{weak labels}; some question mappings are AI-assisted. Recording the label source makes the assignment process repeatable and allows results to be compared by source. A BERT classifier adapted to BRSR language uses these labelled text segments to learn the nine reporting principles. This choice connects the classification task directly to the categories defined by BRSR.

The third design choice makes company comparisons understandable through a visible score breakdown. The Transparency Score combines Completeness, Specificity, Clarity, and mapped Sustainable Development Goal coverage. Each component contributes to the total through a stated weight. Comparisons use the recorded sector so an analyst can examine a company alongside sector peers. Principle probabilities and numerical review flags remain separate outputs, allowing the score, classification, and evidence checks to be examined individually.

GRASP uses a BERT classifier to assign text segments to the nine BRSR principles. Word lists identify vague expressions. Semantic similarity compares the meaning of sentences. Numerical rules create review records containing a claim and the relevant reported figures. The BERT classifier learns from labelled BRSR text segments. A zero-shot MPNet method ranks principles without training on that labelled dataset and provides a comparison during evaluation. The system requires collected reports, a database, and computing resources for training and analysis. After training, the BERT classifier can be reused by the service layer. An application programming interface, or API, allows software to access the stored information and analysis results. Saved results for 1,164 companies show that the workflow has been applied to a large collection of reports.

GRASP aims to help analysts find supporting evidence, group text segments by BRSR principle, and compare company scores with a clear breakdown of the calculations. Results can support the review of one company, comparisons within a sector, and analysis across sectors. Reports available for consecutive years also support year-by-year comparison. Section~\ref{sec:methodology} presents the algorithms, training procedure, and score formulas. Section~\ref{sec:evaluation} reports the measured model results, while Section~\ref{sec:results} presents the company and sector applications.
